% The spoken body of the homily for the resumed Sixth Sunday after the
% Epiphany, said on Sunday 15 November 2026. Continuous preaching, ready to
% read aloud after the Gospel: no heading, no direction to the preacher, no
% citation inside the speech. The loci, the relation to the reviewed
% interpretations, the audience, the word count and the pace estimate stand in
% the terminal note, which no one reads out. The blank lines are the paragraph
% breaks a speaker recovers his place by. The speech tells the two parables;
% raises the difficulty of leaven as a figure of the kingdom and answers it
% with St Ambrose; gives St Augustine's leaven of love and St Bede's reason
% for the three measures, the Church's commandment of love with the whole
% heart, soul and strength; takes St John Chrysostom on the verb "hid" and
% Bede on "until"; meets the slowness of change with Bede on the mustard seed
% and Ambrose on the Church's slow work; brings in the Collect's thought,
% word and deed, and the Epistle's labour of charity with Chrysostom on
% Jason; hears the cry from the depths with Chrysostom; gives one response
% from the Communion; and closes on the Postcommunion and Bede's end of the
% leaven's work. The honorific is tied to the name so that "St" never ends a
% line.
\raggedcolumns

The Lord told the crowds two short parables from a boat on the Sea of Galilee,
and the Gospel gives them one after the other. A man sows a mustard seed in
his field. It is the least of all seeds, and it grows into a tree, so that the
birds come and dwell in its branches. A woman takes leaven and hides it in
three measures of meal, a large batch, until the whole is leavened. Matthew
says that Jesus spoke to the crowds only in parables, as the psalm had said:
``I will open my mouth in parables, I will utter things hidden from the
foundation of the world.'' He does not explain these two. When the disciples
came to Jesus in the house afterwards, they asked him only about the cockle
sown among the wheat. St~John Chrysostom supposes that they let the mustard
seed and the leaven pass as being plainer.

The second parable is less plain than it looks, because elsewhere in the New
Testament leaven stands for corruption. ``Take heed and beware of the leaven
of the Pharisees and Sadducees,'' the Lord tells his disciples, and St~Paul
tells the Corinthians to feast ``not with the old leaven, nor with the leaven
of malice and wickedness.'' Yet here the kingdom of heaven is like leaven.
St~Ambrose answered from those very warnings: they show, he says, that
leaven stands for teaching. ``The leaven of the cockle is one thing,'' he
adds, ``the leaven of the wheat another.'' The leaven of the parable, for
Ambrose, is the teaching of the Church.

St~Augustine gave another answer, in a few words: the leaven is love, which
warms and rouses. St~Bede took up Augustine's words and gave a reason for the
three measures. The woman, he says, is the Church, and she hides the leaven of
love in three measures of meal because she commands us to love the Lord with
the whole heart, the whole soul and the whole strength. So the three measures
are the heart, the soul and the strength that the great commandment asks for
whole. Here Ambrose's answer and Augustine's come together. The teaching is the
commandment, and what the commandment hides in us is love, in the whole
person, with no part left out.

Then there is the verb. Chrysostom, who took the leaven for the apostles mixed
among the nations, saw that the Lord did not say the woman put the leaven in
the meal. He said that she hid it. Leaven is not set beside the meal but mixed
into it, and ``though it be buried, yet is not destroyed, but by little and
little transmutes all into its own condition.'' Hidden leaven cannot be
watched at work. For a time the batch looks like nothing but meal.

And the parable has an ``until.'' She hid it ``until the whole was leavened.''
Bede says that this is well said, because love hidden in our mind must go on
growing until it has changed the whole mind into its own perfection. Then the
soul can love nothing, do nothing and remember nothing but the love of its
maker.

Many of us have heard these parables for years. We have come to Mass and to
confession all that time, and much in us seems unchanged. It is easy to
conclude that nothing is happening. Neither parable is about speed: the seed
has to grow, and the leaven has to work its way through the meal. Bede says of
the mustard seed that the preaching of the gospel grows also in the mind of
each believer, because no one is made perfect all at once. But the leaven was
hidden in all three measures, and a part of us kept back from love stays as it
was.

The Collect at the beginning of this Mass asked for that whole. It asked that,
ever fixing our thoughts on reasonable things, we may both in word and in deed
do what is pleasing to God. Where the English says ``fixing our thoughts,''
the Latin speaks of meditating, and of meditating always. Ambrose says that
the Church's leavening of body, soul and spirit is not hasty and not
haphazard, but brought about by long meditation, so that the three may be one.
Our constant thought, the meditating the Collect asks for, is our part in that
work. But love can stay in our thoughts and never reach our speech, or reach
our speech and stop short of our hands. So the Collect goes on from thought
and asks that what pleases God be done in word and in deed.

St~Paul gave thanks for love of that kind in the church at Thessalonica. He
had preached in their synagogue for three sabbaths, and some believed. Then a
mob attacked the house of Jason, where Paul was staying, and the brethren sent
Paul away by night. Writing to them afterwards, he remembers ``the work of
your faith and labour and charity.'' Paul wrote in Greek, and in the Greek
that Chrysostom read the words are ``the labour of love.'' Chrysostom asks
what labour there is in loving. ``Merely to love is no labor at all. But to
love genuinely is great labor,'' when a thousand things are stirred up that
would draw us from love and we hold out against them all. His example is
Jason. The mob could not find Paul, so they dragged Jason before the rulers of
the city for having received him, and Chrysostom asks whether it was a slight
labour to endure such a storm ``when the seed had not yet taken root.'' Jason
gave security and sent Paul away. In Jason love had gone past thought and word
into a deed, and the deed put him in danger.

The Alleluia has already sung the cry from the depths, and the Offertory will
sing it again: ``From the depths I have cried to Thee, O Lord: Lord, hear my
prayer.'' Chrysostom says that prayer from the deepest heart, even before it
receives what it asks, calms the passions, softens anger and drives out envy.
Prayer of that kind works on the one who prays from within, out of sight, as
leaven works in the meal.

Bede also says what the meal is. Grain is ground between two stones before it
is meal. The meal, he says, is ourselves, ground by fear and hope as between
an upper and a lower millstone, so that, as the Apostle says, we being many
may be one bread, one body in Christ. At Communion we who come to the altar
partake of one bread, and the Mass gives the Lord's promise: ``whatsoever you
ask when you pray, believe that you shall receive, and it shall be done to
you.'' So when you come to Communion today, ask for this: that the love of God
hidden in you may reach the whole of you, your thoughts, your words and your
deeds; and believe that you shall receive it.

After Communion the priest will ask that we who have been fed with celestial
delights may ever hunger after those things by which we truly live. People who
have just been fed ask to go on hungering, because the work is not finished in
any of them. Bede says where it will be finished. Love, he says, is begun here
and perfected there, where God is all in all and warms everyone with one and
the same fire of his love. Until then, the slowness you find in yourselves is
not the leaven failing. It is leaven hidden in three measures of meal, still
at work until the whole is leavened.
